% Created 2022-01-19 Wed 21:11
% Intended LaTeX compiler: pdflatex
\documentclass[11pt]{article}
\usepackage[utf8]{inputenc}
\usepackage[T1]{fontenc}
\usepackage{graphicx}
\usepackage{longtable}
\usepackage{wrapfig}
\usepackage{rotating}
\usepackage[normalem]{ulem}
\usepackage{amsmath}
\usepackage{amssymb}
\usepackage{capt-of}
\usepackage{hyperref}
\author{Francesc Rocher}
\date{\today}
\title{Day 1 : Sonar Sweep}
\hypersetup{
 pdfauthor={Francesc Rocher},
 pdftitle={Day 1 : Sonar Sweep},
 pdfkeywords={},
 pdfsubject={},
 pdfcreator={Emacs 29.0.50 (Org mode 9.5.2)},
 pdflang={English}}
\begin{document}

\maketitle
\tableofcontents


\section{PART 1}
\label{sec:org2eadb98}
\subsection{Statement}
\label{sec:org9918733}
\begin{html}
<details>
<summary>
  <i>click to expand</i>
</summary>
\end{html}

You're minding your own business on a ship at sea when the overboard alarm goes
off! You rush to see if you can help. Apparently, one of the Elves tripped and
accidentally sent the sleigh keys flying into the ocean!

Before you know it, you're inside a submarine the Elves keep ready for
situations like this. It's covered in Christmas lights (because of course it
is), and it even has an experimental antenna that should be able to track the
keys if you can boost its signal strength high enough; there's a little meter
that indicates the antenna's signal strength by displaying 0-50 \textbf{stars}.

Your instincts tell you that in order to save Christmas, you'll need to get \textbf{all
fifty} stars by December 25th.

Collect stars by solving puzzles. Two puzzles will be made available on each day
in the Advent calendar; the second puzzle is unlocked when you complete the
first. Each puzzle grants one star. Good luck!

As the submarine drops below the surface of the ocean, it automatically performs
a sonar sweep of the nearby sea floor. On a small screen, the sonar sweep report
(your puzzle input) appears: each line is a measurement of the sea floor depth
as the sweep looks further and further away from the submarine.

For example, suppose you had the following report:

\begin{verbatim}
199
200
208
210
200
207
240
269
260
263
\end{verbatim}

This report indicates that, scanning outward from the submarine, the sonar sweep
found depths of 199, 200, 208, 210, and so on.

The first order of business is to figure out how quickly the depth increases,
just so you know what you're dealing with - you never know if the keys will get
carried into deeper water by an ocean current or a fish or something.

To do this, \textbf{count the number of times a depth measurement increases} from the
previous measurement. (There is no measurement before the first measurement.) In
the example above, the changes are as follows:

\begin{verbatim}
199 (N/A - no previous measurement)
200 (increased)
208 (increased)
210 (increased)
200 (decreased)
207 (increased)
240 (increased)
269 (increased)
260 (decreased)
263 (increased)
\end{verbatim}

In this example, there are 7 measurements that are larger than the previous
measurement.

\textbf{How many measurements are larger than the previous measurement?}

\begin{html}
</details>
\end{html}

\subsection{Solution}
\label{sec:orga4f5c03}
The algorithm to solve this problem is so easy that don't need to be explicitly
exposed:

\begin{enumerate}
\item read a measurement from input file
\item compare the current measurement with the previous one
\begin{itemize}
\item if the previous is lesser than the current, then an increment occurred
\end{itemize}
\item the result is the total amount of increments detected
\end{enumerate}

\subsection{Implementation}
\label{sec:org0017eed}
\subsubsection{Declarations}
\label{sec:org653fc7a}
The input is stored in a file consisting in a set of lines with a number in
each. To the read the file we use the generic package \texttt{Ada.Text\_IO.Integer\_IO}
instantiated with the type \texttt{Natural}:

\begin{verbatim}

--  __Packages__
package Measurement_IO is new Integer_IO (Natural);
use Measurement_IO;

\end{verbatim}

Variables used in the process are:

\begin{itemize}
\item \texttt{Input}, the input file
\item \texttt{Measurement} and \texttt{Previous\_Measurement}, to store input values
\item \texttt{Increments}, to count the total number of increments detected
\end{itemize}

\begin{verbatim}

--  __Variables__
Input                : File_Type;
Measurement          : Natural;
Previous_Measurement : Natural := Natural'Last;
Increments           : Natural := 0;

\end{verbatim}

\subsubsection{Process}
\label{sec:orgc575d94}
Simply read all input measurements and compare with previous one. If there have
been an increment, then count it.

\begin{verbatim}

--  __Detect_Increments__
Get (Input, Measurement);
loop
  if Previous_Measurement < Measurement then
     Increments := Increments + 1;
  end if;

  exit when End_Of_File (Input);

  --  prepare next iteration
  Previous_Measurement := Measurement;
  Get (Input, Measurement);
end loop;

\end{verbatim}

\subsubsection{Result}
\label{sec:orgdf09220}
The result prints the number of detected increments.

\begin{verbatim}

--  __Result__
Put_Line ("Answer:" & Increments'Image);

\end{verbatim}

\subsubsection{Compilation unit}
\label{sec:org4726cd3}
The compilation unit for this program is obtained with the composition of the
source code blocks shown here:

\begin{verbatim}
with Ada.Text_IO; use Ada.Text_IO;

procedure Day01_P1 is

   --  __Packages__
   package Measurement_IO is new Integer_IO (Natural);
   use Measurement_IO;


   --  __Variables__
   Input                : File_Type;
   Measurement          : Natural;
   Previous_Measurement : Natural := Natural'Last;
   Increments           : Natural := 0;

begin
   Open (Input, In_File, "__Path__" & "input");

      --  __Detect_Increments__
      Get (Input, Measurement);
      loop
        if Previous_Measurement < Measurement then
           Increments := Increments + 1;
        end if;

        exit when End_Of_File (Input);

        --  prepare next iteration
        Previous_Measurement := Measurement;
        Get (Input, Measurement);
      end loop;

   Close (Input);

   --  __Result__
   Put_Line ("Answer:" & Increments'Image);

end Day01_P1;
\end{verbatim}

\subsubsection{Evaluation}
\label{sec:org8394bbd}

\begin{verbatim}
Answer: 1446
\end{verbatim}

\section{PART 2}
\label{sec:org94a7589}
\subsection{Statement}
\label{sec:orgc799151}
\begin{html}
<details>
<summary>
  <i>click to expand</i>
</summary>
\end{html}

Considering every single measurement isn't as useful as you expected: there's
just too much noise in the data.

Instead, consider sums of a three-measurement sliding window. Again considering
the above example:

\begin{verbatim}
199  A
200  A B
208  A B C
210    B C D
200  E   C D
207  E F   D
240  E F G
269    F G H
260      G H
263        H
\end{verbatim}


Start by comparing the first and second three-measurement windows. The
measurements in the first window are marked A (199, 200, 208); their sum is
199 + 200 + 208 = 607. The second window is marked B (200, 208, 210); its sum
is 618. The sum of measurements in the second window is larger than the sum of
the first, so this first comparison \textbf{increased}.

Your goal now is to count \textbf{the number of times the sum of measurements in this
sliding window increases} from the previous sum. So, compare A with B, then
compare B with C, then C with D, and so on. Stop when there aren't enough
measurements left to create a new three-measurement sum.

In the above example, the sum of each three-measurement window is as follows:

\begin{verbatim}
A: 607 (N/A - no previous sum)
B: 618 (increased)
C: 618 (no change)
D: 617 (decreased)
E: 647 (increased)
F: 716 (increased)
G: 769 (increased)
H: 792 (increased)
\end{verbatim}

In this example, there are 5 sums that are larger than the previous sum.

Consider sums of a three-measurement sliding window. \textbf{How many sums are larger
than the previous sum?}

\begin{html}
</details>
\end{html}

\subsection{Solution}
\label{sec:org6066d15}
The solution is quite similar to part 1, with the different that now the
measurements are composed by the sum of three input measurements, or a
\emph{measurement window}. First measurement window is composed by input measurements
1, 2 and 3. Second measurement window by the input measurements 2, 3 and 4, and
so on. The firs measurement windows is complete when the third measurement input
is read, so no comparisons with previous measurement window can be done until
that point. From third input onward, there will be two measurement windows
completed that can be compared.

Each input measurement must be added in three consecutive measurement windows:

\begin{verbatim}
          ,------ 199  A
window 1 <        200  A B ------.
          `---,-- 208  A B C      > window 2
    window 3 <    210    B C D --'
              `-- 200  E   C D

 Input measurement 'A' must be added to windows 1, 2 and 3
\end{verbatim}

That means that three measurement windows must be updated with each input. But
it is not necessary to keep them all, just the last three.

\subsection{Implementation}
\label{sec:org9f3ec32}
Because three measurement windows need to be cyclically updated during the
process, the implementation in Ada can make use of \texttt{mod} types. In particular,
\texttt{mod 3} (values 0, 1 \& 2) will be required for the window index.

Arithmetic operations with this type are used to select windows. For example,
when \texttt{I = 2} then \texttt{I + 1 = 0} refers to the \emph{next} window that will be updated.
That is, the oldest and complete of the three.

\subsubsection{Declarations}
\label{sec:org15bafaf}
Two types are declared to manage the measurement windows,

\begin{verbatim}

--  __Types__
type Window_Index is mod 3;
type Measurement_Window is array (Window_Index) of Natural;

\end{verbatim}

and the variables to manage them:

\begin{verbatim}

--  __New_Variables__
Input_Window         : Measurement_Window := (others => 0);
I                    : Window_Index;

\end{verbatim}

\subsubsection{Process}
\label{sec:org0dfaa22}

At the beginning there are no three windows, still. Thus, the process starts
reading the first three input measurements and add them to the respective
windows:

\begin{verbatim}

--  __Read_First_Input_Values__
Get (Input, Input_Window (0));  --  1st measurement --> 1st window

Get (Input, Input_Window (1));  --  2nd measurement --> 1st & 2nd windows
Input_Window (0) := Input_Window (0) + Input_Window (1);

Get (Input, Input_Window (2));  --  3rd measurement --> 1st, 2nd & 3rd windows
Input_Window (0) := Input_Window (0) + Input_Window (2);
Input_Window (1) := Input_Window (1) + Input_Window (2);

\end{verbatim}

Once the first window is complete, loop over the rest of the input measurements:

\begin{verbatim}

--  __Process_Rest_Of_Input_Values__
I := 2;  --  current window
loop
  Measurement := Input_Window (I + 1);
  if Previous_Measurement < Measurement then
     Increments := Increments + 1;
  end if;

  exit when End_Of_File (Input);

  --  prepare next itaration
  I := I + 1;
  Get (Input, Input_Window (I));
  Input_Window (I - 1) := Input_Window (I - 1) + Input_Window (I);
  Input_Window (I - 2) := Input_Window (I - 2) + Input_Window (I);
  Previous_Measurement := Measurement;
end loop;

\end{verbatim}

\subsubsection{Result}
\label{sec:orgb2101bd}
The result prints the number of detected increments, as before.

\subsubsection{Compilation unit}
\label{sec:org8c23fbe}
The compilation unit for this program is obtained with the composition of the
source code blocks shown here:

\begin{verbatim}
with Ada.Text_IO; use Ada.Text_IO;

procedure Day01_P2 is

   --  __Types__
   type Window_Index is mod 3;
   type Measurement_Window is array (Window_Index) of Natural;


   --  __Packages__
   package Measurement_IO is new Integer_IO (Natural);
   use Measurement_IO;


   --  __Variables__
   Input                : File_Type;
   Measurement          : Natural;
   Previous_Measurement : Natural := Natural'Last;
   Increments           : Natural := 0;


   --  __New_Variables__
   Input_Window         : Measurement_Window := (others => 0);
   I                    : Window_Index;

begin
   Open (Input, In_File, "__Path__" & "input");

      --  __Read_First_Input_Values__
      Get (Input, Input_Window (0));  --  1st measurement --> 1st window

      Get (Input, Input_Window (1));  --  2nd measurement --> 1st & 2nd windows
      Input_Window (0) := Input_Window (0) + Input_Window (1);

      Get (Input, Input_Window (2));  --  3rd measurement --> 1st, 2nd & 3rd windows
      Input_Window (0) := Input_Window (0) + Input_Window (2);
      Input_Window (1) := Input_Window (1) + Input_Window (2);


      --  __Process_Rest_Of_Input_Values__
      I := 2;  --  current window
      loop
        Measurement := Input_Window (I + 1);
        if Previous_Measurement < Measurement then
           Increments := Increments + 1;
        end if;

        exit when End_Of_File (Input);

        --  prepare next itaration
        I := I + 1;
        Get (Input, Input_Window (I));
        Input_Window (I - 1) := Input_Window (I - 1) + Input_Window (I);
        Input_Window (I - 2) := Input_Window (I - 2) + Input_Window (I);
        Previous_Measurement := Measurement;
      end loop;

   Close (Input);

   --  __Result__
   Put_Line ("Answer:" & Increments'Image);

end Day01_P2;
\end{verbatim}

\subsubsection{Evaluation}
\label{sec:org4c13111}

\begin{verbatim}
Answer: 1486
\end{verbatim}

\section{REFERENCES}
\label{sec:org1833ee6}

\begin{itemize}
\item source: \href{https://adventofcode.com/2021/day/1}{Advent of Code 2021, day 1}
\item tangled source code files:
\begin{itemize}
\item \url{src/day01\_p1.adb}
\item \url{src/day01\_p2.adb}
\end{itemize}
\item input file: \url{input}
\item Ada/SPARK support for Org-Babel: \href{https://github.com/rocher/ob-ada-spark}{ob-ada-spark}
\end{itemize}
\end{document}